\begin{apartats}

\apartat[1,25]{0,75}
Calcula les integrals indefinides $\displaystyle\int\frac{3}{x-1}\,dx$ i
$\displaystyle\int\frac{2}{(x+1)^2}\,dx$.

\begin{solucio}
$\displaystyle\int\frac{3}{x-1}\,dx=3\ln|x-1|+k$. Escrivint $\frac2{(x+1)^2}=2(x+1)^{-2}$,
\[
\int\frac{2}{(x+1)^2}\,dx=2\,\frac{(x+1)^{-1}}{-1}+k=-\frac{2}{x+1}+k.
\]
\end{solucio}

\begin{tria}{barrow-racionals}
\itemtria{barrow}{1}{1,25}
Calcula $\displaystyle\int_2^4\frac{3}{x-1}\,dx$ i $\displaystyle\int_0^1\frac{2}{(x+1)^2}\,dx$.

\begin{solucio}
$\displaystyle\int_2^4\frac{3}{x-1}\,dx=\bigl[3\ln|x-1|\bigr]_2^4=3\ln3-3\ln1=3\ln3$.\\
$\displaystyle\int_0^1\frac{2}{(x+1)^2}\,dx=\left[-\frac{2}{x+1}\right]_0^1=-1-(-2)=1$.
\end{solucio}

\itemtria{limit-desconegut}{1}{1,25}
Troba el valor de $a>0$ perquè $\displaystyle\int_0^a(2x+1)\,dx=6$.

\begin{solucio}
$\displaystyle\int_0^a(2x+1)\,dx=\bigl[x^2+x\bigr]_0^a=a^2+a$. Cal $a^2+a=6$, és a dir, $a^2+a-6=0$, amb
solucions $a=2$ i $a=-3$. Com que $a>0$, $a=2$.
\end{solucio}
\end{tria}

\begin{nomesllarg}
\apartat{0,75}
Calcula $\displaystyle\int_1^4\frac{3}{\sqrt x}\,dx$.

\begin{solucio}
$\frac3{\sqrt x}=3x^{-1/2}$, amb primitiva $3\cdot\frac{x^{1/2}}{1/2}=6\sqrt x$:
$\displaystyle\int_1^4\frac{3}{\sqrt x}\,dx=\bigl[6\sqrt x\bigr]_1^4=12-6=6$.
\end{solucio}
\end{nomesllarg}

\end{apartats}
